\documentclass[12pt]{article}

\usepackage{graphicx}
\usepackage{url}
\usepackage{hyperref}

\renewcommand{\topfraction}{1.0}
\renewcommand{\bottomfraction}{1.0}
\renewcommand{\textfraction}{0.0}
\renewcommand{\floatpagefraction}{1.0}
\renewcommand{\dbltopfraction}{1.0}


\begin{document}

\begin{titlepage}

\begin{center}
{\large STAT 5810/6810 Introduction to Statistical Computing} \\[4cm]

{\LARGE \bf Instructions for Homework  2} \\[1cm]
by \\[0.5cm]
{\bf J\"urgen Symanzik} \\[2.5cm]
{\bf Date:} \today \\[2cm]

UTAH STATE UNIVERSITY \\[0.5cm]
Logan, UT \\[0.5cm]
Fall 2012 \\[0.5cm]
\end{center}

\thispagestyle{empty}
\vfill
\end{titlepage}


\newpage 


\pagenumbering{roman}

\tableofcontents


\newpage


%\listoftables
%\addcontentsline{toc}{section}{List of Tables}
%
%\newpage
%
%\listoffigures
%\addcontentsline{toc}{section}{List of Figures}
%
%\newpage 


\pagenumbering{arabic}


\section{\LaTeX\ Installation and Usage} 

\subsection{\LaTeX\ Installation}

\LaTeX\ has been installed on many computers at Utah State University (USU). If you need to install
\LaTeX\ on your own computer, google for it. For PCs, a frequently used version of \LaTeX\ is
MikTeX which can be obtained from \url{http://miktex.org/}.
Follow the download and installation instructios.


\subsection{\LaTeX\ Usage}

There exist many support tools that simplify the usage of \LaTeX. 
I use {\it TeXworks} from \url{http://www.tug.org/texworks/}.
Depending on your \LaTeX\ installation, {\it TeXworks} may be installed
automatically as part of the  \LaTeX\ installation.
If not, follow the download and installation instructions.
Other tools that simplify the use and translation of \LaTeX\ documents
also exists. It is up to you to explore what is available. Talk to some
of the other graduate students who use \LaTeX\ for their MS and PhD work
and let them explain which support tools (if any) they use for their work with \LaTeX.

Most of these tools allow you to translate a \LaTeX\ document (a file with
extension {\tt .tex}) via the click of one (or a few) buttons.

If you can't decide which support tool to install, you will have to 
translate your \LaTeX\ files manually. Change your working directory
to the directory where the main {\tt .tex} file is located.
Then type \verb|pdflatex filename| (assuming your main file is
called {\tt filename.tex}). This will create a file called 
 {\tt filename.pdf} plus several support files. These files
will be used in future steps to create the table of context,
the reference list, an index section, and so on. 

Just to be on the safe side that everything gets properly updated, 
you should run \verb|pdflatex filename| three times before you
finally submit a dcoument --- or otherwise, it may happen that
a last minute shift of a section heading to a new page may not
be reflected in the final table of contents.

If you installed a support tool, this tool almost certainly will
know how many times it needs to repeat the \verb|pdflatex filename|
sequence (and any other steps that  may be needed).


\subsection{Images in \LaTeX}

\subsubsection{pdf, jpg, and png Format}

\LaTeX\ supports most major types of images, but those cannot be mixed
in arbitrary ways. {\tt pdflatex}, as described above, allows you to
incorporate images in pdf, jpg, and png format. 


\subsubsection{ps and eps Format}

If you happen to have images in ps or eps format, then you have
to translate your document via  \verb|latex filename|.
This will create a file called 
{\tt filename.dvi} plus several support files. You can view this
file directly (check the help pages of your installation),
or you can translate it via  \verb|dvips filename| or
 \verb|dvi2ps filename| into a ps document (again, check the help
pages of your installation). You can view this ps file
directly with a postscript viewer or translate it further into pdf
with Adobe Acrobat. In any case, working with ps \& eps
images is more complicated, but such images usually are of a
higher quality than jpg images in particular.


\subsubsection{tiff Format}

\LaTeX\ does not support tiff images, although many publishers request
tiff images. In such a case, you  need to create a working copy of
your image (e.g., in jpg format) and use that for your local 
\LaTeX\ translation. Once  you have submitted your \LaTeX\ source files
to the publisher, they will replace the working copy of your image
with the high--resolution tiff image on their side at the final
production stage of a book chapter or journal article.


\subsubsection{Inclusion of Images}

Below is an example that shows how to include a pdf image.
\LaTeX\ allows you to determine where to place a figure
(here at the {\bf b}ottom of the page), adjust the width (or height)
of a figure (here, 99\% of the maximum textwidth), and later reference this figure via
\verb|Figure~\ref{fig_CCmap}|:

\newpage 

{\small
\begin{verbatim}
\begin{figure}[b]
\includegraphics[width=0.99\textwidth]{fig_CCmap.pdf}
\caption{\label{fig_CCmap}
CCmaps, based on data from the USDA--NASS web page,
related to soybean production in the United States.
The plot shows the dependent variable production [top slider] that is 
conditioned on the two independent variables acreage [bottom slider] 
and yield [right slider].
}
\end{figure}
\end{verbatim}
}

Whenever you include figures (and tables), create meaningful  captions!
Also, keep in mind that each figure must be referenced at least
once in the text. Figures (and tables) {\bf must} be numbered in
the order they are referenced in the text. You can't refer to
Fig.~2 before you  refer to Fig.~1. In such a case, you must
change the order of your figures in your document.


\section{Homework 2}

In most cases, we start with a \LaTeX\ template. There exists templates for
MS reports, MS theses, and PhD dissertations at USU. Check with other graduate
students for good templates. Also, many journals, major publishers, or organizers
for conferences provide \LaTeX\ templates for their publications. Authors
basically are asked to keep the overall formatting in place and
just make changes to the template, e.g., filling in author names,
title, abstract, and, of course, the actual text.

For Homework 2, you have to do the same, i.e., start with this template,
fill in your name, and your answers, but keep the overall formatting, title page, and so on.
During the next few months, we will extend this template with additional
\LaTeX\ features.


\newpage


In particular, for this homework, incorporate the answers in your .R file as follows:

{\small
\begin{verbatim}
Your answers to Homework 2 should be incorporated in a 
verbatim enviroment (see the .tex source file for details).

See w h a t this      environment does
1
 2
   3
     4
       5
End (of verbatim).
\end{verbatim}
}

Just as a comparison, below are the previous seven lines, now outside a verbatim environment: \\[1cm]

 See w h a t this      environment does
1
 2
   3
     4
       5
End (of verbatim). \\[1cm]

\noindent
{\bf \Large All you need to submit for this homework is the final resulting pdf file.}


\section{Getting Help}

There exist numerous sources for help with \LaTeX:

\begin{enumerate}
\item The \LaTeX\ Wikibook at \url{http://en.wikibooks.org/wiki/LaTeX}.

\item The Wikipedia page of \LaTeX\ at \url{http://en.wikipedia.org/wiki/LaTeX}
with many excellent references (online and books).

\item Information from the TeX Users Group (TUG) at
\url{http://www.tug.org/begin.html}.
\end{enumerate}

As mentioned above, many of the other graduate students at USU are using \LaTeX.
Help may be available at the next desk. Otherwise, ask some of the other
students in class, use Canvas to ask a question, or directly contact  me!

{\bf Finally, do not forget to google!} I usually include the word {\tt latex}
in my google query when I search for a particular \LaTeX\ command. \\[2cm]

\noindent
{\Huge Good luck!!!}


\end{document}

